\section{Реализация и результаты}
\label{sec:work}

Линейная рециркуляционная сеть выполняет автоассоциативное отображение вектора блока через узкий скрытый слой. Прямое распространение сжимает вектор, обратное~\mdash  восстанавливает его:
\begin{equation}
	Y = X W,\qquad X' = Y W',
\end{equation}
где $W$ имеет размер $n \times p$, $W'$~\mdash  размер $p \times n$, $p < n$.

Схема преобразования показана на рисунке~\ref{fig:schema}.

\begin{figure}[H]
	\centering
	\includegraphics[width=0.95\textwidth]{schema.png}
	\caption{Схема линейной рециркуляционной сети}
	\label{fig:schema}
\end{figure}

Обучение выполняется по правилу коррекции ошибки. Для каждого эталона вычисляется $\varepsilon = X' - X$, после чего обновляются веса. В реализованном варианте коэффициент обучения адаптивный:
\begin{equation}
	\alpha = \frac{1}{\|X\|^2 + \|Y\|^2}.
\end{equation}
После обновления выполняется нормирование столбцов матрицы $W$ и строк матрицы $W'$ к единичной длине.

Суммарная ошибка эпохи
\begin{equation}
	E = \frac{1}{2}\sum_{k=1}^{Q}\|X'_k - X_k\|^2.
\end{equation}
Обучение продолжается до достижения порога $E \le \varepsilon$ либо до исчерпания лимита эпох.

Коэффициент сжатия рассчитывается по формуле методических указаний:
\begin{equation}
	Z(Q,n,p) = \frac{Q \cdot n}{(n + Q) \cdot p}.
\end{equation}
Требуется $Z > 1$.

Программа реализована на языке Python с использованием библиотек NumPy, Pillow и Matplotlib. Структура модулей:
\begin{itemize}
	\item \texttt{image\_processing.py}~\mdash  загрузка BMP, нарезка блоков, нормализация и денормализация, расчёт $Z$;
	\item \texttt{network.py}~\mdash  модель сети, обучение, сжатие и восстановление;
	\item \texttt{main.py}~\mdash  интерактивный ввод параметров и сохранение результатов.
\end{itemize}

Ключевые фрагменты обучения приведены в листингах~\ref{lst:train-sample} и~\ref{lst:alpha}.

\begin{lstlisting}[caption={Шаг обучения на одном блоке},label={lst:train-sample}]
def _train_sample(self, X: np.ndarray) -> float:
    Y = X @ self.W
    Xp = Y @ self.Wp
    err = Xp - X
    alpha = self._alpha_for_sample(X, Y)
    dWp = -alpha * np.outer(Y, err)
    dW = -alpha * np.outer(X, err @ self.Wp.T)
    self.Wp = self.Wp + dWp
    self.W = self.W + dW
    if NORMALIZE_WEIGHTS:
        self._normalize_weights()
    return float(np.dot(err, err))
\end{lstlisting}

\begin{lstlisting}[caption={Адаптивный коэффициент обучения},label={lst:alpha}]
def _alpha_for_sample(self, X: np.ndarray, Y: np.ndarray) -> float:
    if not ADAPTIVE:
        return self.alpha
    denom = float(np.dot(X, X) + np.dot(Y, Y))
    return 1.0 / max(denom, 1e-12)
\end{lstlisting}

Эксперимент выполнен на тестовом цветном изображении $256 \times 256$. Параметры: $r = m = 8$, $p = 32$, максимум $200$ эпох, $\varepsilon = 0.01$. Получено $Q = 1024$ блоков, $n = 192$, коэффициент сжатия $Z \approx 5{,}05 > 1$.

Пример работы консольного интерфейса показан на рисунке~\ref{fig:console}.

\begin{figure}[H]
	\centering
	\includegraphics[width=0.92\textwidth]{console.png}
	\caption{Пример запуска программы}
	\label{fig:console}
\end{figure}

Исходный фрагмент изображения и результат восстановления приведены на рисунках~\ref{fig:original} и~\ref{fig:restored}. График суммарной ошибки по эпохам~\mdash  на рисунке~\ref{fig:error}.

\begin{figure}[H]
	\centering
	\includegraphics[width=0.55\textwidth]{original.png}
	\caption{Исходное изображение (обрезанная область)}
	\label{fig:original}
\end{figure}

\begin{figure}[H]
	\centering
	\includegraphics[width=0.55\textwidth]{restored.png}
	\caption{Восстановленное изображение}
	\label{fig:restored}
\end{figure}

\begin{figure}[H]
	\centering
	\includegraphics[width=0.92\textwidth]{error_curve.png}
	\caption{Зависимость суммарной ошибки $E$ от номера эпохи}
	\label{fig:error}
\end{figure}

По результатам прогона: $E \approx 150{,}4$, средняя ошибка по пикселям $\mathrm{MSE} \approx 29{,}5$. Ошибка быстро убывает в начале обучения и выходит на плато, что соответствует сходимости модели при выбранном $p$. Визуально восстановленное изображение сохраняет крупные цветовые области и контуры фигур при допустимых искажениях, обусловленных сжатием с потерями ($p < n$).
